% from ega4-18-fr.tex
\oldpage[IV-4]{206}

si l'on pose $\mathbf f'=(f_1,\ldots,f_{n-1})$, la suite $\mathbf f'$ est $M'$-régulière et l'on a, en vertu de $c)$ appliqué à $\mathbf f'$, une suite exacte
\[
  0\longrightarrow M'/(\textstyle\sum_{i=1}^{n-1}f_iM')\longrightarrow M/(\textstyle\sum_{i=1}^{n-1}f_iM)\longrightarrow M''/(\textstyle\sum_{i=1}^{n-1}f_iM'')\longrightarrow0.
\]
Par hypothèse, $f_n$ est $(M/(\sum_{i=1}^{n-1}f_iM))$-régulier et $(M''/(\sum_{i=1}^{n-1}f_iM''))$-régulier, donc le même raisonnement montre d'une part que $f_n$ est $(M'/(\sum_{i=1}^{n-1}f_iM'))$-régulier, et d'autre part, en vertu de (3.4.1.4), que la suite
\[
  0\longrightarrow M'/(\textstyle\sum_{i=1}^{n}f_iM')\longrightarrow M/(\textstyle\sum_{i=1}^{n}f_iM)\longrightarrow M''/(\textstyle\sum_{i=1}^{n}f_iM'')\longrightarrow0
\]
est exacte, d'où $c)$.

Montrons enfin que $a)$ entraîne $b)$. Comme la suite $\mathbf f$ est alors M-régulière et $M'$-régulière, donc M-quasi-régulière et $M'$-quasi-régulière, on a des isomorphismes canoniques
\[
  \operatorname{gr}_{\mathfrak J}^*(M')\xrightarrow{\sim}\operatorname{gr}_{\mathfrak J}^0(M')[T_1,\ldots,T_n],
  \qquad
  \operatorname{gr}_{\mathfrak J}^*(M)\xrightarrow{\sim}\operatorname{gr}_{\mathfrak J}^0(M)[T_1,\ldots,T_n]
\]
($0$, 15.1.7), et comme on a vu ci-dessus que $\operatorname{gr}_{\mathfrak J}^0(M')\to\operatorname{gr}_{\mathfrak J}^0(M)$ est injectif, il en est de même de $\operatorname{gr}_{\mathfrak J}^*(M')\to\operatorname{gr}_{\mathfrak J}^*(M)$.

On notera encore que l'équivalence des conditions $a)$, $b)$, $c)$ de (19.5.3) est valable en particulier lorsque $A$ est noethérien, $M''$ un $A$-module de type fini et que les $f_i$ appartiennent au radical de $A$.

\begin{env}[19.5.4]
\label{IV.19.5.4}
Dans la suite de ce numéro, nous conservons les notations $A$, $\mathbf f$, $M$, $\mathfrak J$ de (19.5.1), et nous supposons de plus $M$ muni d'une filtration décroissante $(M_k)$ formée de sous-$A$-modules, telle que $M_0=M$. Rappelons ($0$, 15.1.5) que l'on définit alors sur $M$ une seconde filtration décroissante formée des sous-modules
\[
  M_k'=M_k+\mathfrak JM_{k-1}+\cdots+\mathfrak J^{k-1}M_1+\mathfrak J^kM_0
\]
et que, si $\operatorname{gr}_\bullet(M)$ et $\operatorname{gr}_\bullet'(M)$ sont les $A$-modules gradués associés aux filtrations $(M_k)$ et $(M_k')$ respectivement, on définit un homomorphisme gradué surjectif de degré $0$ ($0$, 15.1.5.2)
\[
  \psi_M:(\operatorname{gr}_\bullet(M)\otimes_A(A/\mathfrak J))[T_1,\ldots,T_n]\longrightarrow\operatorname{gr}_\bullet'(M).
\]
Rappelons aussi que l'on a l'homomorphisme canonique surjectif ($0$, 15.1.1.1)
\[
  \varphi_{\operatorname{gr}_\bullet(M)}:(\operatorname{gr}_\bullet(M)\otimes_A(A/\mathfrak J))[T_1,\ldots,T_n]\longrightarrow\operatorname{gr}_{\mathfrak J}^*(\operatorname{gr}_\bullet(M)).
\]
\end{env}

\begin{theorem}[19.5.5]
\label{IV.19.5.5}
Avec les notations de (19.5.4), considérons les propriétés suivantes:
\begin{enumerate}
  \item[a)] La suite $\mathbf f$ est $\operatorname{gr}_\bullet(M)$-régulière.
  \item[b)] L'homomorphisme canonique $\psi_M$ est bijectif.
  \item[c)] L'homomorphisme canonique $\varphi_{\operatorname{gr}_\bullet(M)}$ est bijectif (autrement dit ($0$, 15.1.7), la suite $\mathbf f$ est $\operatorname{gr}_\bullet(M)$-quasi-régulière).
\end{enumerate}

